%!TEX root=thesis.tex
Cognitive impairment is a growing health issue causing serious challenges to patients, societies, and healthcare systems worldwide. The majority of cognitive impairment cases are undetected due to the lack of access to affordable diagnostic tools, especially in low- and middle-income countries. Current predictive models suffer from limitations that hinder their use in clinical settings, including \hl{the limited assessment of model calibration and decision curve analysis} and the need for model interpretability. \hl{Cognitive impairment develops gradually from neuropathological changes that begin years before a clinical diagnosis is reached, and in clinical practice the cognitive status is determined by following recognized diagnostic frameworks. Accessible and routinely collected data can capture the early signs of this gradual process and approximate the clinical diagnosis.} This study aims to address this gap by developing cost-effective and clinically reliable machine learning models that can accurately predict cognitive impairment early using accessible data. The study proposes a methodological framework designed based on a hierarchical classification approach comprising five distinct binary classification tasks. These tasks consist of classifying the current cognitive status, identifying the primary etiology of cognitive impairment, and predicting the risk of cognitive impairment up to five years in the future. The models are trained using demographic, health history, and behavioral and functional assessments data from the National Alzheimer's Coordinating Center's Uniform Data Set. The methodology utilizes ensemble feature selection and hybrid data resampling methods guided by class overlap measures to systematically preprocess training data. Diverse traditional machine learning models are evaluated using a novel model selection algorithm to identify the most cost-effective classifiers. The selected models are validated across NACC research centers and on an external dataset from the Alzheimer's Disease Neuroimaging Initiative. The models are also calibrated and assessed using calibration curve and decision curve analysis. Finally, the methodology incorporates a multilevel model explanation approach to interpret the predictions of the models and generate personalized counterfactual explanations that identify modifiable risk factors for prevention. The results demonstrated that the application of the proposed data preprocessing methods successfully reduced the complexity of the training data and improved the classification performance. Consequently, logistic regression models achieved comparable performance to complex and non-linear models with major improvements in computational efficiency. These models generalized well to independent datasets, achieving an area under the curve of 0.825 and 0.778 to predict current and future cognitive impairment, respectively. In addition, the calibration curve and decision curve analysis revealed accurate estimates with a calibration error below 5\% and a high net benefit across a wide range of probability thresholds, confirming the clinical utility of the developed models. Furthermore, personalized counterfactual explanations were generated for more than 86\% of the cognitively impaired subjects. These findings demonstrate the feasibility of developing cost-effective, interpretable, and clinically reliable machine learning models to early predict the risk of cognitive impairment. \hl{However, these models were developed and validated on retrospective cohort data, and they are intended to support clinical decision-making rather than replace it. These constraints should be considered when the findings are interpreted.} This study contributes to the body of research by providing a data-driven methodological framework to serve as a foundation for reliable, accessible, and equitable cognitive impairment screening models.